% =========================================================
% T-MASTER — FULL TRAVERSAL (T0–T8)
% Hyper-Condensed DSLO v0.7 Corridor
% =========================================================

\section{T-MASTER Traversal Corridor}

% ---------------------------------------------------------
\subsection{T0 Origin Manifold}
The origin manifold defines the substrate-level geometry from which all other
manifolds derive. It provides origin curvature, identity seeds, legality
anchors, collapse surfaces, and recovery envelopes.

\paragraph{Origin Manifold Diagram}
This visual shows the vertical traversal of the Origin Manifold: curvature seed
field, identity seed region, legality anchor plane, collapse surface, and
recovery envelope. Each node represents a foundational geometric layer in the
origin substrate.


\begin{center}
\begin{tikzpicture}[
    node distance=1.2cm,
    every node/.style={rectangle, draw, align=center, minimum width=6cm, rounded corners}
]

\node (curv) {\textbf{Curvature Seed Field}};
\node (id) [below=of curv] {\textbf{Identity Seed Region}\\(Initial identity wells and boundaries)};
\node (leg) [below=of id] {\textbf{Legality Anchor Plane}\\(Lawful origin constraints)};
\node (col) [below=of leg] {\textbf{Collapse Surface}\\(Origin collapse geometry)};
\node (rec) [below=of col] {\textbf{Recovery Envelope}\\(Origin recovery attractor)};

\draw[-{Latex}] (curv) -- (id);
\draw[-{Latex}] (id) -- (leg);
\draw[-{Latex}] (leg) -- (col);
\draw[-{Latex}] (col) -- (rec);

\end{tikzpicture}

\vspace{8pt}
\textit{(Replace with full manifold diagram in v0.7 expansion)}
\end{center}





% ---------------------------------------------------------
\subsection{T1 Unified Substrate Manifold}
The unified manifold defines global curvature, identity boundaries, legality
masks, collapse geometry, recovery geometry, supervisor geometry, projection
geometry, and continuity constraints.

\paragraph{Unified Substrate Manifold Diagram}
This visual depicts the unified substrate manifold stack: global curvature,
identity boundaries, legality masks, collapse geometry, recovery geometry,
supervisor geometry, projection geometry, and continuity constraints. It forms
the backbone of all derived manifolds.


\begin{center}
\begin{tikzpicture}[
    node distance=1.2cm,
    every node/.style={rectangle, draw, align=center, minimum width=6cm, rounded corners}
]

\node (curv) {\textbf{Global Curvature Field}};
\node (id) [below=of curv] {\textbf{Identity Boundary Geometry}};
\node (leg) [below=of id] {\textbf{Legality Mask Suite}};
\node (col) [below=of leg] {\textbf{Collapse Geometry}};
\node (rec) [below=of col] {\textbf{Recovery Geometry}};
\node (sup) [below=of rec] {\textbf{Supervisor Geometry}};
\node (proj) [below=of sup] {\textbf{Projection Geometry}};
\node (cont) [below=of proj] {\textbf{Continuity Constraints}};

\draw[-{Latex}] (curv) -- (id);
\draw[-{Latex}] (id) -- (leg);
\draw[-{Latex}] (leg) -- (col);
\draw[-{Latex}] (col) -- (rec);
\draw[-{Latex}] (rec) -- (sup);
\draw[-{Latex}] (sup) -- (proj);
\draw[-{Latex}] (proj) -- (cont);

\end{tikzpicture}
\end{center}


% ---------------------------------------------------------
\subsection{T2 Derived Manifold Suite}
The derived suite defines six lawful projections: Agency, Teleology,
Deployment, Execution, Federation, and Omega.

\paragraph{Derived Manifold Suite Diagram}
This visual shows the six derived manifolds branching from the unified substrate:
Agency, Teleology, Deployment, Execution, Federation, and Omega. It illustrates
the lawful projection structure of DSLO’s manifold suite.


\begin{center}
\begin{tikzpicture}[
    node distance=1.2cm,
    every node/.style={rectangle, draw, align=center, minimum width=4cm, rounded corners}
]

\node (root) {\textbf{Unified Substrate Manifold}};

\node (agency) [below left=1.2cm and 2cm of root] {\textbf{Agency Manifold}};
\node (tele)   [below=1.2cm of root] {\textbf{Teleology Manifold}};
\node (deploy) [below right=1.2cm and 2cm of root] {\textbf{Deployment Manifold}};

\node (exec) [below=1.2cm of tele] {\textbf{Execution Manifold}};
\node (fed)  [below left=1.2cm and 2cm of exec] {\textbf{Federation Manifold}};
\node (omega)[below right=1.2cm and 2cm of exec] {\textbf{Omega Manifold}};

\draw[-{Latex}] (root) -- (agency);
\draw[-{Latex}] (root) -- (tele);
\draw[-{Latex}] (root) -- (deploy);

\draw[-{Latex}] (tele) -- (exec);
\draw[-{Latex}] (exec) -- (fed);
\draw[-{Latex}] (exec) -- (omega);

\end{tikzpicture}
\end{center}


% ---------------------------------------------------------
\subsection{T3 Relational Geometry}
Relational geometry defines cross-system interaction: relational primitives,
fields, invariants, drift geometry, legality geometry, and human-machine
relational dynamics.

\paragraph{Relational Geometry Diagram}
This visual outlines relational geometry: relational primitives, fields,
invariants, drift geometry, legality geometry, and human–machine relational
dynamics. It represents the lawful interaction structure between systems.


\begin{center}
\begin{tikzpicture}[
    node distance=1.2cm,
    every node/.style={rectangle, draw, align=center, minimum width=6cm, rounded corners}
]

\node (prim) {\textbf{Relational Primitives}};
\node (fields) [below=of prim] {\textbf{Relational Fields}};
\node (inv) [below=of fields] {\textbf{Relational Invariants}};
\node (drift) [below=of inv] {\textbf{Drift Geometry}};
\node (leg) [below=of drift] {\textbf{Relational Legality Geometry}};
\node (hm) [below=of leg] {\textbf{Human-Machine Relational Dynamics}};

\draw[-{Latex}] (prim) -- (fields);
\draw[-{Latex}] (fields) -- (inv);
\draw[-{Latex}] (inv) -- (drift);
\draw[-{Latex}] (drift) -- (leg);
\draw[-{Latex}] (leg) -- (hm);

\end{tikzpicture}
\end{center}


% ---------------------------------------------------------
\subsection{T4 Operators of Reality}
Operators define lawful transformations across all manifold layers: fold,
invert, lift, bind, collapse-trajectory, recovery-window, curvature-redirect,
identity-boundary, and legality-preservation.

\paragraph{Operators of Reality Diagram}
This visual lists the operator stack: fold, invert, lift, bind, collapse-
trajectory, recovery-window, curvature-redirect, identity-boundary, and
legality-preservation. It represents lawful transformations across manifold
layers.


\begin{center}
\begin{tikzpicture}[
    node distance=1.2cm,
    every node/.style={rectangle, draw, align=center, minimum width=5cm, rounded corners}
]

\node (fold) {\textbf{Fold Operator}};
\node (invert) [below=of fold] {\textbf{Invert Operator}};
\node (lift) [below=of invert] {\textbf{Lift Operator}};
\node (bind) [below=of lift] {\textbf{Bind Operator}};
\node (ct) [below=of bind] {\textbf{Collapse-Trajectory Operator}};
\node (rw) [below=of ct] {\textbf{Recovery-Window Operator}};
\node (cr) [below=of rw] {\textbf{Curvature-Redirect Operator}};
\node (ib) [below=of cr] {\textbf{Identity-Boundary Operator}};
\node (lp) [below=of ib] {\textbf{Legality-Preservation Operator}};

\draw[-{Latex}] (fold) -- (invert);
\draw[-{Latex}] (invert) -- (lift);
\draw[-{Latex}] (lift) -- (bind);
\draw[-{Latex}] (bind) -- (ct);
\draw[-{Latex}] (ct) -- (rw);
\draw[-{Latex}] (rw) -- (cr);
\draw[-{Latex}] (cr) -- (ib);
\draw[-{Latex}] (ib) -- (lp);

\end{tikzpicture}
\end{center}


% ---------------------------------------------------------
\subsection{T5 Thermodynamic Geometry}
Thermodynamic geometry defines pressure fields, pressure flow, drift geometry,
collapse geometry, recovery geometry, human drift, machine drift, and dual
drift comparison.

\paragraph{Thermodynamic Geometry Diagram}
This visual shows thermodynamic geometry: pressure fields, pressure flow, drift
geometry, collapse geometry, recovery geometry, human drift, machine drift, and
dual drift comparison. It represents DSLO’s pressure–drift–collapse–recovery
loop.


\begin{center}
\begin{tikzpicture}[
    node distance=1.2cm,
    every node/.style={rectangle, draw, align=center, minimum width=6cm, rounded corners}
]

\node (press) {\textbf{Pressure Fields}};
\node (flow) [below=of press] {\textbf{Pressure Flow}};
\node (drift) [below=of flow] {\textbf{Drift Geometry}};
\node (col) [below=of drift] {\textbf{Collapse Geometry}};
\node (rec) [below=of col] {\textbf{Recovery Geometry}};
\node (hd) [below=of rec] {\textbf{Human Drift}};
\node (md) [below=of hd] {\textbf{Machine Drift}};
\node (dual) [below=of md] {\textbf{Dual Drift Comparison}};

\draw[-{Latex}] (press) -- (flow);
\draw[-{Latex}] (flow) -- (drift);
\draw[-{Latex}] (drift) -- (col);
\draw[-{Latex}] (col) -- (rec);
\draw[-{Latex}] (rec) -- (hd);
\draw[-{Latex}] (hd) -- (md);
\draw[-{Latex}] (md) -- (dual);

\end{tikzpicture}
\end{center}


% ---------------------------------------------------------
\subsection{T6 DSUP Runtime Geometry}
Runtime geometry defines state representation, lawful update cycles, constrained
evolution, context-integrated updates, drift-coherence correction, halt
conditions, completion cycles, and meaning-condition coupling.

\paragraph{DSUP Runtime Geometry Diagram}
This visual depicts runtime geometry: state representation, lawful update
cycles, constrained evolution, context-integrated updates, drift-coherence
correction, halt conditions, completion cycles, and meaning-condition coupling.
It represents DSUP’s lawful runtime evolution.


\begin{center}
\begin{tikzpicture}[
    node distance=1.2cm,
    every node/.style={rectangle, draw, align=center, minimum width=6cm, rounded corners}
]

\node (state) {\textbf{State Representation}};
\node (update) [below=of state] {\textbf{Lawful Update Cycles}};
\node (evo) [below=of update] {\textbf{Constrained Evolution}};
\node (ctx) [below=of evo] {\textbf{Context-Integrated Updates}};
\node (dcc) [below=of ctx] {\textbf{Drift-Coherence Correction}};
\node (halt) [below=of dcc] {\textbf{Halt Conditions}};
\node (comp) [below=of halt] {\textbf{Completion Cycles}};
\node (mcc) [below=of comp] {\textbf{Meaning-Condition Coupling}};

\draw[-{Latex}] (state) -- (update);
\draw[-{Latex}] (update) -- (evo);
\draw[-{Latex}] (evo) -- (ctx);
\draw[-{Latex}] (ctx) -- (dcc);
\draw[-{Latex}] (dcc) -- (halt);
\draw[-{Latex}] (halt) -- (comp);
\draw[-{Latex}] (comp) -- (mcc);

\end{tikzpicture}
\end{center}


% ---------------------------------------------------------
\subsection{T7 Context Window Geometry}
Context geometry defines context windows, boundaries, drift, contraction
operators, legality masks, context-state mapping, and drift flags.

\paragraph{Context Window Geometry Diagram}
This visual outlines context window geometry: windows, boundaries, drift,
contraction operators, legality masks, context-state mapping, and drift flags.
It represents lawful context traversal and drift control.


\begin{center}
\begin{tikzpicture}[
    node distance=1.2cm,
    every node/.style={rectangle, draw, align=center, minimum width=6cm, rounded corners}
]

\node (win) {\textbf{Context Windows}};
\node (bound) [below=of win] {\textbf{Context Boundaries}};
\node (drift) [below=of bound] {\textbf{Context Drift}};
\node (contract) [below=of drift] {\textbf{Contraction Operators}};
\node (mask) [below=of contract] {\textbf{Legality Masks}};
\node (map) [below=of mask] {\textbf{Context-State Mapping}};
\node (flag) [below=of map] {\textbf{Drift Flags}};

\draw[-{Latex}] (win) -- (bound);
\draw[-{Latex}] (bound) -- (drift);
\draw[-{Latex}] (drift) -- (contract);
\draw[-{Latex}] (contract) -- (mask);
\draw[-{Latex}] (mask) -- (map);
\draw[-{Latex}] (map) -- (flag);

\end{tikzpicture}
\end{center}


% ---------------------------------------------------------
\subsection{T8 Simulation Geometry}
Simulation geometry defines IRSM, simulation integrity conditions, static
anchoring, topological anchors, four-plane legality checks, projection planes,
mapping planes, invariant validation planes, and CLCP propagation.

\paragraph{Simulation Geometry Diagram}
This visual shows simulation geometry: IRSM, simulation integrity conditions,
static anchoring, topological anchors, and the four-plane legality corridor
(text projection, geometric mapping, invariant validation, CLCP). It represents
lawful simulation traversal across manifold layers.


\begin{center}
\begin{tikzpicture}[
    node distance=1.2cm,
    every node/.style={rectangle, draw, align=center, minimum width=6cm, rounded corners}
]

% --- Vertical stack ---
\node (irsm) {\textbf{IRSM (Invariant-Restricted Simulation Mode)}};
\node (int) [below=of irsm] {\textbf{Simulation Integrity Conditions}};
\node (sa) [below=of int] {\textbf{Static Anchoring}};
\node (ta) [below=of sa] {\textbf{Topological Anchors}};
\node (fp) [below=of ta] {\textbf{Four-Plane Legality Check}};

% --- Centered horizontal matrix for the three planes ---
\node (planes) [below=1.2cm of fp] {
    \begin{tikzpicture}[
        every node/.style={rectangle, draw, align=center, minimum width=4.5cm, rounded corners},
        node distance=1.2cm
    ]
        \node (tpp) {\textbf{Text Projection Plane}};
        \node (gmp) [right=1.2cm of tpp] {\textbf{Geometric Mapping Plane}};
        \node (ivp) [right=1.2cm of gmp] {\textbf{Invariant Validation Plane}};
    \end{tikzpicture}
};

% --- CLCP plane centered under GMP ---
\node (clcp) [below=1.2cm of planes] {\textbf{CLCP Plane}};

% --- Vertical arrows ---
\draw[-{Latex}] (irsm) -- (int);
\draw[-{Latex}] (int) -- (sa);
\draw[-{Latex}] (sa) -- (ta);
\draw[-{Latex}] (ta) -- (fp);

% --- Arrows from FP to the three planes ---
\draw[-{Latex}] (fp) -- (planes);

% --- Arrow from GMP to CLCP ---
\draw[-{Latex}] (planes) -- (clcp);

\end{tikzpicture}
\end{center}



% =========================================================
% END OF T-MASTER TRAVERSAL
% =========================================================
